\documentclass[10pt]{beamer}

%\usetheme{ENSLyon}

\usetheme{AnnArbor}

\title[CNMAC 2021]{\Large \textsc{Software HCprot:\\ Pré processamento de Instâncias Moleculares do Protein Data Bank}}
\author[G. Philippi]{{Guilherme Philippi\vspace{-0.3cm}}}
\institute[]{Departamento de Matemática\\ Universidade Federal de Santa Catarina, Blumenau \\ \vspace{0.2cm} Desenvolvido em conjunto com Felipe Fidalgo e Emerson V. Castelani\vspace{-0.2cm}}
\date[15 Setembro, 2021]{ {\small XL Congresso Nacional de Matemática Aplicada e Computacional} \\ {\scriptsize Minissimpósio 10 - Geometria de Distâncias e Álgrebras Geométricas\\ 15 de setembro de 2021}}
\setbeamersize{text margin left=5mm}
\setbeamersize{text margin right=5mm}

\setbeamertemplate{navigation symbols}{}
%\usecolortheme{ENSLyon_greener}
\usecolortheme{whale}

\setbeamercolor{frametitle}{bg=gray}
\setbeamertemplate{bibliography item}{\insertbiblabel}

%PACKAGES -------------------------
\usepackage{etex}
\usepackage[utf8]{inputenc}
\usepackage[brazilian]{babel}
\usepackage{enumerate}
\usepackage{amsmath}
\usepackage{amssymb}
\usepackage{amsthm}
\usepackage{amscd}
\usepackage{amsfonts}
\usepackage{multicol}
\usepackage{multirow}
\usepackage{array}
\usepackage{color}
\usepackage{graphicx}
\usepackage{tikz}
\usepackage{tikz-qtree}
\usepackage{wrapfig}
\usepackage{3dplot}
\usepackage{pgf}
\usepackage{tkz-euclide}
\usepackage{algorithmic}
\usepackage{algorithm}
\usepackage{xparse}
\usepackage{subfigure}

%LIBRARIES-TIKZ ------------------------------------------

\usetikzlibrary{shadows,trees}
\usetikzlibrary{decorations.pathmorphing}
\usetikzlibrary{decorations.markings}
\usetikzlibrary{positioning}
\usetikzlibrary{chains,matrix,scopes}
\usetikzlibrary{arrows}

%DEFINITIONS ----------------------------------------------------

\def\centerarc[#1](#2)(#3:#4:#5)% Syntax: [draw options] (center) (initial angle:final angle:radius)
{ \draw[#1] ($(#2)+({#5*cos(#3)},{#5*sin(#3)})$) arc(#3:#4:#5); }
\def\xx{\mathbf{x}}
\def\ii{\mathbf{i}}
\def\jj{\mathbf{j}}
\def\kk{\mathbf{k}}
\def\tt{\mathbf{t}}
\def\ee{\mathbf{e}}
\def\qq{\mathbf{q}}
\def\pp{\mathbf{p}}
\def\vv{\mathbf{v}}
\def\rr{\mathbf{r}}
\def\vzero{\mathbf{0}}
\def\qset{\mathbb{H}}
\def\xx{\mathbf{x}}


%\AtBeginSection[]
%{
%	\begin{frame}<beamer>
%		\frametitle{Inicio da seção \thesection}
%		\tableofcontents[currentsection]
%	\end{frame}
%}


%NEW THEOREMS ------------------------------------------

\theoremstyle{plain}
\newtheorem{teorema}{Teorema}[section]
\newtheorem{lema}{Lema}[section]
\newtheorem{proposicao}{Proposição}[section]
\newtheorem{corolario}{Corolário}[section]

\theoremstyle{definition}
\newtheorem{definicao}{Definição}[section]
\newtheorem{observacao}{Observação}[section]
\newtheorem{exemplo}{Exemplo}[section]

\newenvironment{solucao}
{\renewcommand\qedsymbol{$\triangle$}\begin{proof}[Solução]}{\end{proof}}



\begin{document}
	
	%FACE
	\begin{frame}
		
		\titlepage
		
		\vspace{-0.7cm}
		\begin{flushleft}
			\includegraphics[scale=0.12]{logo.png}
		\end{flushleft}
		
		\vspace{-2cm}
		\begin{flushright}
			\includegraphics[scale=0.024]{logo_ufsc.png}
		\end{flushright}
	\end{frame}
	
	%Indice
	\begin{frame}
		\tableofcontents 
	\end{frame}
	
	\section{Introdução}
	
	\begin{frame}
		\frametitle{\normalsize Introduction } 
		{
			
		}	
	\end{frame}
	
	\section{Referências}
	%Slide refe
	\begin{frame}		
		
		{\footnotesize
		\bibliographystyle{unsrt}
		\bibliography{references}}
	\end{frame}
	
	%Slide End
	\begin{frame}
		\begin{center}
			\vspace{1.5cm}
			Obrigado!\\
			\hspace{-4.5cm}
			\includegraphics[scale=0.2]{logo.png}
			
			\vspace{-2.7cm}
			\hspace{5.5cm}
			\includegraphics[scale=0.038]{logo_ufsc.png}
			
			\vspace{0.5cm}
			Contato: g.philippi\@grad.ufsc.br\\ UFSC - Blumenau
		\end{center}
	\end{frame}
	
\end{document}